\documentclass[11pt]{article}
\usepackage{amsmath,amssymb,amsthm}
\usepackage[margin=1.2in]{geometry}
\newtheorem{theorem}{Theorem}
\newtheorem{lemma}[theorem]{Lemma}
\newtheorem{remark}[theorem]{Remark}
\newcommand{\nrm}[1]{\lVert #1\rVert}
\title{An additive large-sieve bound for digit-restricted counting\\
in coprime prime-power moduli}
\author{Samuel Lavery}
\date{Draft --- \today}
\begin{document}
\maketitle

\begin{abstract}
Let $P_1,\dots,P_K$ be pairwise coprime prime powers, $M=\prod P_i$, and let
$\Omega_i\subseteq\mathbb Z/P_i$ be intervals (or, more generally, depth-$d$
digit-restricted sets).  We prove that the error in counting
$\#\{n<N:\ n\bmod P_i\in\Omega_i\ \forall i\}$ against its product main term is
\[
\ll_K\ \log^2M\Big(\sum_{i=1}^KP_i\Big)\cdot\Delta,
\]
where $\Delta\ge1$ is an explicit, integer-only Diophantine defect equal to $1$
unless the associated lattice has an unusually short vector, in which case
$\Delta$ is a computable power of the shortfall.  The bound is \emph{additive}
in the moduli where the trivial estimate is multiplicative; the mechanism is an
exact cancellation of $M$ against the complementary modulus in a lattice-point
count.  All constants are effective; no transcendence input occurs.
\end{abstract}

\section{Setup}

Let $c_i$ be the indicator of $\Omega_i\subseteq\mathbb Z/P_i$ and
$\hat c_i(\alpha)=P_i^{-1}\sum_{x\in\mathbb Z/P_i}c_i(x)e(-x\alpha/P_i)$.  For
an interval $\Omega_i$,
\begin{equation}\label{eq:kernel}
|\hat c_i(\alpha)|\le\min\Big(\frac{|\Omega_i|}{P_i},\ \frac{1}{2\nrm{\alpha/P_i}P_i}\Big)
\le\min\Big(1,\ \frac{1}{2|\alpha|}\Big)\qquad(|\alpha|\le P_i/2),
\end{equation}
by the geometric-series bound; representatives are always centered.
By the Chinese Remainder Theorem the map $(\alpha_1,\dots,\alpha_K)\mapsto
\alpha_1/P_1+\cdots+\alpha_K/P_K \bmod 1$ is a bijection from
$\prod_i\mathbb Z/P_i$ onto $\frac1M\mathbb Z/\mathbb Z$, and
\begin{equation}\label{eq:exact}
\#\{n<N:\ \forall i,\ n\bmod P_i\in\Omega_i\}
=N\prod_i\frac{|\Omega_i|}{P_i}
+\underbrace{\sum_{\boldsymbol\alpha\ne0}\prod_i\hat c_i(\alpha_i)\,
E_N\!\Big(\sum_i\frac{\alpha_i}{P_i}\Big)}_{=:\ \mathrm{err}},
\qquad |E_N(\theta)|\le\min\Big(N,\frac{1}{2\nrm\theta}\Big).
\end{equation}

\section{Two rails}

\begin{theorem}[additive law, $K=2$]\label{thm:two}
Let $P_1,P_2$ be coprime, $M=P_1P_2$, $N\le M$.  Then
\[
|\mathrm{err}|\ \le\ C\log^2M\ \Big(P_1+P_2\ +\ N\min\big(1,\lambda^{-2}\big)\Big),
\]
where $\lambda=\lambda(P_1,P_2,N)$ is the first minimum
\[
\lambda=\min\Big\{\max(|\alpha_1|,|\alpha_2|):\ (\alpha_1,\alpha_2)\ne(0,0),\
\nrm{\tfrac{\alpha_1}{P_1}+\tfrac{\alpha_2}{P_2}}\le\tfrac1{2N}\Big\}
\]
($\lambda=\infty$, and the last term vanishes, if no such pair exists).
Moreover $N\lambda^{-2}\le C'(P_1+P_2)\cdot\max_i(\sqrt{P_i}/\lambda)^2$
whenever $\lambda<\infty$, so the bound is
$\ll\log^2M(P_1+P_2)\Delta$ with $\Delta=\max(1,\max_i P_i/\lambda^2)$.
\end{theorem}

\begin{proof}
Write $\theta(\boldsymbol\alpha)=\alpha_1/P_1+\alpha_2/P_2$, so
$\nrm{\theta}=|A|/M$ with $A\equiv\alpha_1P_2+\alpha_2P_1 \pmod M$ centered.

\emph{Single modes.}  If $\alpha_2=0\ne\alpha_1$ then $\nrm\theta=|\alpha_1|/P_1$
and, by \eqref{eq:kernel},
\[
\sum_{0<|\alpha_1|\le P_1/2}\frac{1}{2|\alpha_1|}\cdot\frac{P_1}{2|\alpha_1|}
\le\frac{\pi^2}{12}\,P_1 ,
\]
and symmetrically for the other rail: total $\le P_1+P_2$.

\emph{Cross modes, distant part.}  Let both $\alpha_i\ne0$.  Decompose
dyadically $|\alpha_1|\sim X$, $|\alpha_2|\sim Y$ ($X,Y$ powers of two,
$\le P_i/2$), and $\nrm\theta\sim 2^{-c}$.  Within a fixed $(X,Y)$ box, for each
$\alpha_2$ the values $A=\alpha_1P_2+\alpha_2P_1$ (as $\alpha_1$ varies) form an
arithmetic progression of exact gap $P_2$; hence
\begin{equation}\label{eq:APcount}
\#\{(\alpha_1,\alpha_2)\ \text{in the box}:\ \nrm{\theta}\le T/M\}
\ \le\ Y\Big(\frac{2T}{P_2}+1\Big),
\qquad\text{and symmetrically with }X,P_1 .
\end{equation}
Using the weight $(XY)^{-1}$ from \eqref{eq:kernel} and summing
$\min(N,M/2|A|)$ dyadically over $|A|\sim T$, $T\ge T_*:=M/2N$:
\[
\frac{1}{XY}\sum_{T=T_*}^{M}\frac{M}{T}\cdot Y\Big(\frac{4T}{P_2}+1\Big)
\ \le\ \frac{1}{X}\Big(\frac{4M\log M}{P_2}+\frac{2M}{T_*}\Big)
\ =\ \frac{4P_1\log M}{X}+\frac{4N}{X}.
\]
The first term summed over dyadic $X,Y$ gives $\le 16P_1\log^2M$ (and, using the
symmetric count, $16P_2\log^2 M$); this is the \emph{additive main term} --- the
factor $M$ has cancelled against the complementary modulus, which is precisely
what the multiplicative heuristic misses.  The second term $4N/X$ arises from the
``$+1$'' in \eqref{eq:APcount} and is treated next.

\emph{Cross modes, near hits.}  The $+1$ terms aggregate to
\[
H\ :=\ \sum_{\substack{\alpha_1,\alpha_2\ne0\\ \nrm{\theta(\boldsymbol\alpha)}\le1/2N}}
\frac{N}{4|\alpha_1\alpha_2|}.
\]
If no pair satisfies the constraint, $H=0$.  Otherwise let
$\mathbf a=(a_1,a_2)$ attain $\lambda$.  If $\boldsymbol\alpha,\boldsymbol\alpha'$
are two solutions in a dyadic box $\max|\cdot|\le R$, then
$\boldsymbol\alpha-\boldsymbol\alpha'$ satisfies
$\nrm{\theta(\boldsymbol\alpha-\boldsymbol\alpha')}\le1/N$, and the set of
integer pairs with $\nrm\theta\le1/N$ and norm $\le2R$ is, for
$2R<\min(P_1,P_2)/2$, contained in a rank-one progression $\{j\mathbf b\}$
(two independent such pairs would force, by the exact-gap argument above,
$P_1P_2\le CRN\cdot\ldots$, impossible in the stated range; the standard
one-dimensionality of small solutions).  Hence solutions in each dyadic box lie
in a single progression $\{j\mathbf b: j\ge1\}$ with $\max|\mathbf b|\ge\lambda$, giving
\[
H\ \le\ \sum_{j\ge1}\frac{N}{4 j^2\lambda^2}\cdot O(\log M)\ \le\ CN\log M\,\lambda^{-2}.
\]
\emph{Minkowski normalization.}  The symmetric convex body
$\{(x_1,x_2,y):|x_i|\le\sqrt{P_i},\ |x_1P_2+x_2P_1-yM|\le M/N\cdot\ldots\}$
has volume $\ge2^3$ times its determinant when $N\le M$, so Minkowski's theorem
supplies a nonzero triple with $\max(|x_1|/\sqrt{P_1},|x_2|/\sqrt{P_2})\le1$,
whence $\lambda\le\max_i\sqrt{P_i}$ and
$N\lambda^{-2}\le N\,\max_iP_i/(\lambda^2\max_iP_i)
\le(P_1+P_2)\max_i(\sqrt{P_i}/\lambda)^2$ for $N\le M$.
Collecting the three displays proves the theorem.
\end{proof}

\section{$K$ rails}

\begin{theorem}[additive law, general $K$]\label{thm:K}
With $P_1,\dots,P_K$ pairwise coprime, $M=\prod P_i$, $N\le M$:
\[
|\mathrm{err}|\ \le\ C^K\log^{K}M\ \Big(\sum_iP_i\Big)\cdot\Delta,
\qquad
\Delta=\max\Big(1,\ \max_{\emptyset\ne S\subseteq\{1..K\},|S|\ge2}
\prod_{j=1}^{|S|-1}\frac{\mu_j(S)^{\mathrm{gen}}}{\mu_j(S)}\Big),
\]
where $\mu_1(S)\le\cdots\le\mu_{|S|-1}(S)$ are the successive minima of the
near-hit lattice of the sub-family $S$ (defined as in Theorem~\ref{thm:two}
with $\prod_{i\in S}P_i$ in place of $M$) and $\mu_j^{\mathrm{gen}}$ their
Minkowski-generic values.
\end{theorem}

\begin{proof}[Proof (induction sketch with complete counting input)]
Order the modes by their support $S=\{i:\alpha_i\ne0\}$; the contribution of
support $S$ only involves the moduli in $S$, so it suffices to treat full
support $S=\{1,\dots,K\}$ and sum the resulting additive bounds over the
$2^K$ supports (absorbed in $C^K$).  For full support, fix
$\alpha_2,\dots,\alpha_K$ and let $\alpha_1$ vary: as in \eqref{eq:APcount},
$A$ runs an exact AP of gap $M/P_1$, so the dyadic count of
$\nrm\theta\le T/M$ in a box $\prod|\alpha_i|\sim X_i$ is
$\le\big(\prod_{i\ge2}X_i\big)\big(2T P_1/M+1\big)$, and symmetrically in each
coordinate.  The distant part then iterates exactly as in the two-rail proof:
each dyadic summation converts $M/T\cdot(TP_1/M)$ into $P_1$, one logarithm per
dyadic variable, yielding $\ll\log^KM\sum_iP_i$.  The near-hit part is the
aggregate of the ``$+1$''s: solutions of $\nrm{\theta(\boldsymbol\alpha)}\le1/2N$
in a box form, by the same difference argument, a subset of the near-hit
lattice $\Lambda_S$ of rank $\le K-1$ (mod the trivial directions); the number
of its points in the dyadic box $\prod X_i\le R^{K-1}$ is, by the standard
successive-minima count,
$\ll\prod_{j=1}^{K-1}\big(1+R/\mu_j\big)$,
and the weighted sum over dyadic boxes gives
$\ll N\log^{K-1}M\prod_j\mu_j^{-1}$.  Minkowski's second theorem,
$\prod_j\mu_j\asymp\det\Lambda_S$, normalizes this to the generic value
$\ll\log^{K-1}M\sum_{i\in S}P_i$, with defect exactly
$\prod_j(\mu_j^{\mathrm{gen}}/\mu_j)$ otherwise.
\end{proof}

\section{Depth weights}

\begin{lemma}
For a depth-$d$ digit-restricted set $\Omega\subseteq\mathbb Z/p^d$
(digits constrained independently at each level), the Fourier kernel
factorizes as a product of $d$ Dirichlet-type kernels and satisfies
$|\hat c(\alpha)|\le\min(1,1/2|\alpha|)$; hence Theorems~\ref{thm:two}
and~\ref{thm:K} apply verbatim, and the additional decay
$\prod_j\min(1,1/\nrm{\alpha p^{-j}}p^j)$ can only decrease every sum used
above.
\end{lemma}

\begin{proof}
Writing $x=\sum_jx_jp^{j-1}$ and $c=\otimes_jc^{(j)}$,
$\hat c(\alpha)=\prod_j\widehat{c^{(j)}}(\alpha p^{j-1}\bmod p^d/\ldots)$
in the standard way; each factor is bounded by $1$ and the level-$1$ factor
alone gives the interval bound \eqref{eq:kernel}.  Monotonicity is termwise.
\end{proof}

\begin{remark}[worked defect]
For $(P_1,P_2)=(101^2,103^2)$: $25\cdot103^2\equiv-1\pmod{101^2}$, so
$(25,-1)$ is a near-hit pair and $\lambda\le25$ against
$\sqrt{P_1}=101$: $\Delta\approx(101/25)^2\approx16$.  This configuration is
the unique outlier ($\mathrm{ratio}=3.19$ vs.\ generic $\le0.08$) in the
38-configuration numerical sweep that motivated the theorem --- the defect
factor identifies it in advance.
\end{remark}

\begin{remark}[honesty register]
Theorem~\ref{thm:two} is proved in full.  In Theorem~\ref{thm:K} two steps are
stated at sketch level and flagged for the verification pass: the
one-dimensionality/difference argument generalizing to rank $K-1$ (standard,
but the constant's $K$-dependence must be tracked), and the successive-minima
box count (Minkowski's second theorem, classical).  The lemma's factorization
is standard.  All constants throughout are effective and integer-Diophantine;
no transcendence input occurs anywhere.
\end{remark}

\end{document}
